\documentclass[11pt,a4paper]{article}
\usepackage[T1]{fontenc}
\usepackage[utf8]{inputenc}
\usepackage{lmodern,amsmath,amssymb,amsthm,booktabs}
\usepackage[margin=27mm]{geometry}
\usepackage[hidelinks]{hyperref}
\hypersetup{pdftitle={Classical Paths with a Finite Action Defect},
pdfauthor={navstokgap research working notes}}
\setlength{\emergencystretch}{3em}
\newtheorem{proposition}{Proposition}
\newcommand{\E}{\mathbb E}
\newcommand{\R}{\mathbb R}
\title{Classical Paths with a Finite Action Defect\\
\large An exact two-regulator experiment}
\author{navstokgap research working notes}
\date{7 September 2026}
\begin{document}
\maketitle
\begin{abstract}
A shrinking Gaussian bridge can converge uniformly to the free classical path
while its excess kinetic action converges to a positive constant. With $d$
internal positions and fluctuation action scale $\kappa$, that constant is
controlled by $d\kappa$: the exact mean excess is $d\kappa/2$.
We derive the partition Hessian, normalization and limiting law, then audit a
two-regulator path-integral proposal. A bounded-speed oscillatory example
isolates the role of derivative control. The calculation turns the search for
an action gap into a concrete selection question: which physical premise
fixes a positive surviving quantity across changes of resolution?
\end{abstract}

\section{The object retained by a limit}
The central object is the \emph{action defect}
$\lim_N(S[q_N]-S[q_{\rm cl}])$ for specified convergent paths and fixed
endpoints. Uniform path convergence and convergence of kinetic action impose
different demands on their derivatives. This distinction is familiar from
oscillations in the calculus of variations \cite[\S\S2.1,3.3]{Rindler2017};
here it gives an exactly soluble mechanics test.

Work on the free line during $[0,T]$, with $T>0$, mass $m>0$ and endpoints
$x,z\in\R$. Positions have length units. The parameter $\kappa>0$ has action
units and defines a positive Gaussian path model. We separately examine a
complex oscillatory integral with action parameter $\epsilon>0$.
The paths integrated over are off-shell comparison paths; the free solution
is $q_{\rm cl}(t)=x+t(z-x)/T$. The results concern these specified models.
For the proposed universal quantum scale, the physical selection principle
and the justification of a phase rule remain the research obligations.

\section{Exact normalization at any partition}
Let $\pi=(0=t_0<\cdots<t_N=T)$, $N\ge2$, $\tau_j=t_j-t_{j-1}$ and
$d=N-1$. The piecewise-linear path with nodes $q_0=x,q_N=z$ has action
\[
 S_\pi(q)=\frac m2\sum_{j=1}^N\frac{(q_j-q_{j-1})^2}{\tau_j}.
\]
Set $\eta_j=q_j-q_{\rm cl}(t_j)$, with $\eta_0=\eta_N=0$. Summation of
the cross term gives
\begin{equation}\label{eq:split}
 S_\pi=S_{\rm cl}+\tfrac12\eta^TA_\pi\eta,
 \qquad S_{\rm cl}=\frac{m(z-x)^2}{2T}.
\end{equation}
The $d\times d$ matrix has entries
\[
 (A_\pi)_{jj}=m(\tau_j^{-1}+\tau_{j+1}^{-1}),\qquad
 (A_\pi)_{j,j+1}=(A_\pi)_{j+1,j}=-m\tau_{j+1}^{-1},
\]
and zeros elsewhere. Its units are mass/time, since its quadratic form acts
on node displacements. The continuous Hessian with the $L^2(dt)$ inner product
instead has eigenvalues with units mass/time squared.

\begin{proposition}[Partition determinant and covariance]\label{prop:partition}
The matrix $A_\pi$ is positive definite and
\begin{equation}\label{eq:detgreen}
 \det A_\pi=\frac{m^dT}{\prod_{j=1}^N\tau_j},\qquad
 (A_\pi^{-1})_{ij}=\frac1m\left(\min(t_i,t_j)-\frac{t_it_j}{T}\right).
\end{equation}
The normalized endpoint bridge density with respect to $d\eta_1\cdots d\eta_d$
is
\begin{equation}\label{eq:density}
 \rho_\pi^\kappa(\eta)=
 \frac{\sqrt{\det A_\pi}}{(2\pi\kappa)^{d/2}}
 \exp\left(-\frac{\eta^TA_\pi\eta}{2\kappa}\right),
 \qquad \operatorname{Cov}(\eta)=\kappa A_\pi^{-1}.
\end{equation}
\end{proposition}
\begin{proof}
The quadratic form is $m\sum_j(\eta_j-\eta_{j-1})^2/\tau_j$; its vanishing
forces every node to zero. For the leading $k\times k$ minor, the tridiagonal
determinant recurrence gives
$D_k=m^k(\tau_1+\cdots+\tau_{k+1})/(\tau_1\cdots\tau_{k+1})$.
This holds at $k=0,1$, and substitution into
$D_k=m(\tau_k^{-1}+\tau_{k+1}^{-1})D_{k-1}
-m^2\tau_k^{-2}D_{k-2}$ proves it inductively.
For each column of the proposed inverse, the piecewise-linear Green function
vanishes at the endpoints, has constant slope on either side of its node,
and a slope decrease of $1/m$ there. Multiplication by $A_\pi$ therefore
gives the corresponding coordinate vector. Diagonalizing the positive matrix
then gives the Gaussian integral and covariance in \eqref{eq:density}.
\end{proof}

There are $N$ transition densities and $d=N-1$ integrated positions:
\begin{equation}\label{eq:kernel}
 \left(\prod_{j=1}^N\sqrt{\frac{m}{2\pi\kappa\tau_j}}\right)
 \int_{\R^d}e^{-S_\pi(q)/\kappa}\,dq
 =\sqrt{\frac{m}{2\pi\kappa T}}e^{-S_{\rm cl}/\kappa}.
\end{equation}
Thus the determinant restores exactly the free transition kernel at duration
$T$. Dividing the product by that kernel yields \eqref{eq:density}.
This is the Brownian transition/bridge construction
\cite[\S\S3.4,4.1]{PitmanYor2018}, with variance parameter $\kappa/m$.

\section{A classical path limit with residual action}
The number of modes and their individual scale enter the action through their
product. This gives the promised two-regulator calculation.

\begin{proposition}[Exact defect law and joint limit]\label{prop:defect}
Under \eqref{eq:density}, the excess $\Delta S_\pi=S_\pi-S_{\rm cl}$ satisfies
\begin{equation}\label{eq:chi}
 \frac{2\Delta S_\pi}{\kappa}\sim\chi_d^2,\qquad
 \E\Delta S_\pi=\frac{d\kappa}{2},\qquad
 \operatorname{Var}(\Delta S_\pi)=\frac{d\kappa^2}{2}.
\end{equation}
Let $N\to\infty$, $d=N-1$, and choose partitions with maximum interval tending
to zero. If $\kappa_N\to0$ and $d\kappa_N\to\ell\in[0,\infty)$, their
polygonal bridges admit a coupling for which
\[
 q_N\longrightarrow q_{\rm cl}\quad\hbox{uniformly almost surely},\qquad
 \Delta S_{\pi_N}\longrightarrow\ell/2\quad\hbox{in }L^2.
\]
The quantity $\ell$ has action units.
\end{proposition}
\begin{proof}
Whitening gives $Z=A_\pi^{1/2}\eta/\sqrt\kappa$ with independent standard
normal coordinates. Equation \eqref{eq:split} becomes
$\Delta S_\pi=\kappa\sum_{j=1}^d Z_j^2/2$. The normal moments
$\E Z_j^2=1$, $\E Z_j^4=3$ prove \eqref{eq:chi} and imply
\[
 \E(\Delta S_{\pi_N}-\ell/2)^2
 =d\kappa_N^2/2+(d\kappa_N-\ell)^2/4\longrightarrow0.
\]
For the path limit, take one continuous standard Brownian bridge
$B_t^T=B_t-(t/T)B_T$ on $[0,T]$ and let $I_\pi$ denote linear interpolation
at the partition nodes. Its covariance is $\min(s,t)-st/T$
\cite[p.~6, (4)--(5)]{PitmanYor2018}. Consequently
$q_N=q_{\rm cl}+\sqrt{\kappa_N/m}\,I_{\pi_N}B^T$ has the required law, and
\[
 \|q_N-q_{\rm cl}\|_\infty
 \le\sqrt{\kappa_N/m}\,\|B^T\|_\infty\longrightarrow0
\]
almost surely by continuity on a compact interval. The action moment identity
holds under this coupling as well.
\end{proof}

For equal steps $\tau=T/N$ and $\kappa_N=b\tau$, where $b>0$ has energy
units, $\ell=bT$ and the residual action is $bT/2$. More generally,
$d\kappa_N\to0$ gives zero residual action, whereas $d\kappa_N\to\infty$
gives divergence in probability: the relative variance is $2/d$.
At fixed $d,\kappa>0$, the support of $\Delta S_\pi$ is $[0,\infty)$.
Thus the positive residual is a selected limiting value, rather than a hard
lower bound on the positive actions of admissible finite paths.

\section{A bounded-speed control example}\label{sec:speed}
A speed bound alone still permits action defects. For the Newtonian action
on $H^1([0,T])$ with fixed endpoints, choose $c>|z-x|/T$ and
$0<v_*<c-|z-x|/T$. The smooth off-shell paths
\begin{equation}\label{eq:sine}
 q_n(t)=q_{\rm cl}(t)+\frac{v_*T}{2\pi n}\sin(2\pi nt/T)
\end{equation}
converge uniformly to $q_{\rm cl}$ and satisfy $|\dot q_n|<c$, while
\begin{equation}\label{eq:sineaction}
 S[q_n]-S[q_{\rm cl}]=\frac{mv_*^2T}{4}>0.
\end{equation}
Indeed the cross term integrates to zero and
$\int_0^T\cos^2(2\pi nt/T)dt=T/2$. Their acceleration amplitude grows as
$2\pi n v_*/T$. This isolates derivative and force control as the next
dynamical test. Here $c$ is an imposed kinematic bound on the Newtonian model;
a relativistic Lagrangian is a separate choice.

The example is an elementary oscillation mechanism of the kind organized by
Young measures \cite[\S3.3]{Rindler2017}. Its coefficient varies continuously
with $v_*$. Strong convergence of velocities in $L^2$ would instead make this
quadratic action converge to the action of the limiting path.

\section{Quadratic stationary phase and critical-point weights}\label{sec:phase}
For one nondegenerate quadratic stationary point, a normalized oscillatory
amplitude gives a precise version of Rivero's finite-dimensional construction
\cite[p.~1, (2)--(4)]{RiveroFeynman1998}. Let
$F(q)=F_*+(q-q_*)^TA(q-q_*)/2$ with $A$ real symmetric positive definite on
$\R^d$, and $O\in\mathcal S(\R^d)$. Here $F$ has action units and $A$
has action/length squared units. Define
\[
 I_\epsilon[O]=(2\pi\epsilon)^{-d/2}
 \int_{\R^d}e^{iF(q)/\epsilon}O(q)\,dq.
\]
Use $\widehat O(k)=(2\pi)^{-d/2}\int e^{-ik\cdot q}O(q)\,dq$.
The Gaussian Fourier identity \cite[p.~464, (14.7)]{GuilleminSternberg2010} gives
\begin{align*}
 e^{-iF_*/\epsilon}I_\epsilon[O]
 &=\frac{e^{i\pi d/4}}{\sqrt{\det A}}(2\pi)^{-d/2}
 \int_{\R^d}e^{ik\cdot q_*}
 e^{-i\epsilon k^TA^{-1}k/2}\widehat O(k)\,dk\\
 &\longrightarrow\frac{e^{i\pi d/4}}{\sqrt{\det A}}O(q_*).
\end{align*}
One can obtain the identity by inserting Gaussian damping in the quadratic
Fourier integral and removing it in tempered distributions. The limit follows
from $\widehat O\in L^1$ and dominated convergence. In particular,
\begin{equation}\label{eq:criticaldelta}
 \lim_{\epsilon\downarrow0}|I_\epsilon[O]|^2
 =\frac{|O(q_*)|^2}{\det A}
 =\langle\delta^{(d)}(\nabla F),|O|^2\rangle.
\end{equation}
The last equality is the change of variables $p=A(q-q_*)$.
Equation \eqref{eq:criticaldelta} is a scalar limit for each fixed test
function $O$; the distribution on its right acts linearly on $|O|^2$.
The amplitude carries a square-root determinant weight and a removable global
phase; its squared modulus gives the gradient-delta weight. Rivero uses
$\epsilon^{-d/2}$, so the corresponding squared limit acquires $(2\pi)^d$
under the present Fourier convention. This is a delta supported on the
critical point, not a derivative of a delta. The argument fixes $d$; joint
refinement requires control of the dimension-dependent weights and observables.

\section{The two regulators and an explicit scaling map}\label{sec:map}
Rivero's equation (5) contains
\[
 \exp\left(\frac{i\tau}{\epsilon}
 \sum_j L^\tau\!\left(q_j,\frac{q_{j+1}-q_j}{\tau}\right)\right),
\]
and (7) takes the two regulators proportional. Under the literal physical
identification $L^\tau=mv^2/2$ with fixed $m$ and fixed physical coordinates,
the exponent is $iS_\pi/\epsilon$. Therefore $[\epsilon]=\text{action}$,
$[\tau]=\text{time}$, and $\epsilon=b\tau$ assigns energy units to $b$.
The positive Gaussian counterpart is precisely the finite-defect scaling in
Proposition~\ref{prop:defect}, with its path width tending to zero.

The source also permits changing $L$ and the coordinate scale, and leaves its
map (6) to be specified. Here is an exact free-model choice. Fix reference
values $m_R>0$, $\kappa_R>0$ and set the bare kinetic coefficient to
\begin{equation}\label{eq:bare}
 m_B(\tau)=\frac{\epsilon(\tau)}{\kappa_R}m_R.
\end{equation}
At every partition this yields
\[
 \frac{S_{\pi,B}}{\epsilon}=\frac{S_{\pi,R}}{\kappa_R},\qquad
 \sqrt{\frac{m_B}{2\pi\epsilon\tau}}
 =\sqrt{\frac{m_R}{2\pi\kappa_R\tau}}.
\]
Both the positive Gaussian family and its oscillatory counterpart are therefore
fixed exactly, with the latter using $i$ in the prefactor and the branch
$1/\sqrt{i}=e^{-i\pi/4}$ for positive time. As $\epsilon\to0$, this map sends
$m_B\to0$ while holding the reference mass $m_R$ fixed. In the oscillatory
case the bare phase corresponds to a fixed reference quantum coefficient;
in the positive case the observable bridge variance is
\begin{equation}\label{eq:renormobs}
 \operatorname{Var}(q(s)\mid x,z)=\frac{\kappa_R s(T-s)}{m_RT}.
\end{equation}
Choosing this variance at one interior time fixes the ratio $\kappa_R/m_R$;
a separately fixed reference mass then fixes $\kappa_R$.

Wilson--Kogut \cite[\S12.2, pp.~166--168, Fig.~12.7]{WilsonKogut1974}
provides the design principle: compare cutoff theories at a common physical
correlation length. Their construction specifies a renormalized mass and
compares rescaled trajectories on the surface with dimensionless correlation
length one. Under the topology discussed there, the renormalized mass remains
a free parameter. In our model, \eqref{eq:renormobs} plays the role of the
reference observable and \eqref{eq:bare} implements its preservation. This is
an explicit toy scaling map, with the physical scale supplied by the
renormalization condition.

\section{The next physical selection test}
The surviving information is now explicit: path locations, velocity
oscillations, total excess action and normalized fluctuation width can have
different limits. The next bounded task is to impose a physical dynamical
condition on the oscillations and calculate which of these quantities survives.
Uniform force control, a finite-speed relativistic action, and a prescribed
fluctuation observable supply concrete alternative tests. For a universal
positive action parameter, a selection argument must also relate different
masses, durations and interacting systems, and explain the phase coefficient.

For the companion PDE question, this example motivates tracking the topology
in which derivatives and quadratic quantities pass to a limit. The
Yang--Mills target additionally concerns the physical Hamiltonian spectrum;
the present defect is an action difference in a fixed-endpoint mechanics model.

\bibliographystyle{plain}
\bibliography{references/library}
\end{document}
